\chapter{Event engines and general body simulation}\label{ch:engine}
\section{Definition, conditions, and evolving state}
\class{Simulation} is the public scenario/run entry point. \class{SimulationOptions} is editable configuration; \class{SimulationConditions} holds the resolved physical services, launch values, numerical limits, and listener list for execution. \class{SimulationStatus} is the evolving state: simulation time, position, world location, velocity, quaternion orientation, angular velocity, active flight configuration, motor states, event queue, flight flags, warnings, and current result branch.

The MIT implementation currently constructs \class{ModifiedEventSimulationEngine} directly in \code{Simulation.simulateInternal()}. The retained \class{BasicEventSimulationEngine} remains a distinct engine and is exercised by focused tests. The modified engine selects RK6 or RK4 using \code{NewControlStepListener.useRK6}, whose default is true. An upstream diagram or nominal engine-class property is therefore not sufficient to identify the actual desktop execution path.

\section{Events around continuous integration}
The engine alternates event processing with numerical steps. Events represent discontinuities and milestones: launch, motor ignition, burnout, ejection charge, stage separation, apogee, recovery deployment, ground impact, tumble, abort, and simulation end. A queue orders pending events; listeners can participate in event addition and handling.

For each active branch, the engine selects the appropriate stepper and determines the time available before the next event. The stepper advances within that bound and can expose conditions for new events. Stage separation creates additional statuses/result branches, with the appropriate active stage subsets and inherited trajectory history. The engine processes those branches independently after separation; it is not a contact/constraint solver for interconnected moving bodies.

\diagram{fig:engine-loop}{A simulation branch alternates event handling and integration. Recovery changes the stepper while preserving the continuing branch state and FC run.}{%
\begin{tikzpicture}
\simple{prepare}{0,0}{Prepare conditions\\initialize listeners and status}{5.4cm}
\simple{events}{7,0}{Process due events\\select active stepper}{5.4cm}
\simple{step}{14,0}{Bound and advance step\\evaluate intermediate states}{5.4cm}
\simple{end}{0,-2.4}{Finalize branch / run\\close listeners and outputs}{5.4cm}
\simple{record}{7,-2.4}{Accept state, record data\\post-step listeners}{5.4cm}
\simple{models}{14,-2.4}{Mass, thrust, wind, atmosphere\\aerodynamics and acceleration}{5.4cm}
\draw[use] (prepare) -- (events);
\draw[use] (events) -- (step);
\draw[use] (step) -- (models);
\draw[use] (models) -- (record);
\draw[use] (record) -- (events);
\draw[use] (record) -- node[flowlabel,above]{termination} (end);
\end{tikzpicture}}

\section{What the flight stepper computes}
\class{RK4SimulationStepper} and \class{RK6SimulationStepper} separately implement their Runge--Kutta schemes. They share \class{AbstractSimulationStepper} services, but RK6 is not a subclass of RK4. Both evaluate multiple provisional states within one accepted time interval. The calculations include:
\begin{enumerate}
\item Query atmospheric and wind conditions and transform relative airflow into the body frame.
\item Calculate aerodynamic coefficients, structure/motor mass, current thrust, and gravity.
\item Convert coefficients to dimensional forces and moments; shift pitch/yaw moments to the current CG.
\item Divide force by mass and the supported moment components by effective inertias; rotate acceleration to the world frame and apply gravity/Coriolis terms.
\item Integrate translation, rotation rate, and quaternion orientation; update world position from local displacement.
\end{enumerate}

Conceptually the translation follows \(\dot{\boldsymbol r}=\boldsymbol v\) and \(\dot{\boldsymbol v}=\boldsymbol a\). The angular implementation uses the calculated pitch/yaw/roll moments and longitudinal/rotational inertia values. It should not be described as a general full-tensor, flexible, or coupled multibody dynamics solution. The source's supported moment and frame conventions are the actual contract.

Before liftoff, the code prevents downward motion through the pad. While the rocket remains on the launch rod, acceleration is projected onto the rod direction and angular acceleration is constrained. The MIT initial-roll-rate option prescribes longitudinal spin, including on a tilted rod; it is not generated by a modeled rail friction or spin-up mechanism.

\section{Timestep selection and fork-specific behavior}
The flight steppers contain limits for configured step size, event distance, angular change, rate change, launch-rod travel, and growth relative to the previous step. However, the non-FC branch retains legacy MIT static timestep controls from \class{MidControlStepLauncher}; it can replace the initially selected step with \code{theTimeStep}. This must be considered when interpreting the ordinary simulation-options timestep field or comparing the fork to vanilla behavior.

With an active FC, the engine/steppers instead honor the FC deadline cap and maximum physics step. The FC path avoids floors or overrides that would skip an imminent virtual deadline. Its default cap is 2.5 ms. The ensemble listener supplies an additional step cap based on the thrust-noise interval. Reducing a step improves temporal resolution but does not repair an incorrect physical submodel or an inconsistent frame transform.

\section{Listener lifecycle and accepted measurements}
\class{SimulationListenerHelper} dispatches lifecycle and calculation hooks. Pre-calculation hooks may supply replacements; post-calculation hooks may inspect or modify values according to their contracts. Step and event callbacks have different semantics from intermediate aerodynamic/acceleration callbacks.

Stateful onboard logic must not execute once for every provisional Runge--Kutta evaluation. The FC listener captures a coherent first-stage acceleration/status acquisition and publishes it only after the interval is accepted. It advances due work at explicit virtual deadlines. This preserves the distinction between a physical derivative evaluation and one firmware loop iteration. The \code{preStep} callback prepares an interval; \code{limitStep} is called by the engine/stepper to enforce its boundary.

\section{Descent, termination, and result publication}
Recovery deployment switches to \class{BasicLandingStepper}; tumbling uses \class{BasicTumbleStepper}; grounded behavior has its own stepper. Landing/tumble derive from \class{AbstractEulerStepper}, which updates translation with its specialized drag model and event-aware limits. Those paths do not integrate a new attitude history, so recorded orientation is held rather than inventing a canopy-hanging or tumbling animation.

The engine records warnings and aborts separately from ordinary samples. Maximum duration, ground impact, explicit aborts, exceptions, and cancellation can end work through different paths. Consumers should inspect the recorded events/warnings and worker outcome rather than infer success merely from the existence of a \class{FlightData} object. Listener shutdown closes per-run resources; the FC removes live execution from auxiliary optimum-coast calculations to avoid advancing its clock a second time.

\sourceline{\source{core/src/main/java/info/openrocket/core/document/Simulation.java}{Simulation.simulateInternal}; \source{core/src/main/java/info/openrocket/core/simulation/ModifiedEventSimulationEngine.java}{ModifiedEventSimulationEngine}; \source{core/src/main/java/info/openrocket/core/simulation/SimulationStatus.java}{SimulationStatus}; \source{core/src/main/java/info/openrocket/core/simulation/RK4SimulationStepper.java}{RK4SimulationStepper}; \source{core/src/main/java/info/openrocket/core/simulation/RK6SimulationStepper.java}{RK6SimulationStepper}; \source{core/src/main/java/info/openrocket/core/simulation/AbstractEulerStepper.java}{AbstractEulerStepper}.}

\section{Ordinary launch, powered flight, and coast}
Run initialization selects active stages and motor installations, builds the initial event queue, establishes pad position and orientation, and initializes listeners. Orientation incorporates launch-rod inclination and bearing; the status also selects an initial roll orientation relative to the flight-condition convention. The initial-spin setting is a separate MIT addition.

The launch event can activate configured motors after their ignition delays. Ignition begins a motor's local clock; burnout ends thrust; ejection follows the selected delay where applicable. Liftoff and rod clearance are physical milestones distinct from ignition. Before liftoff the vehicle cannot translate downward through the pad, and while constrained by the rod the solver projects acceleration onto the guide direction.

\class{SimulationStatus} computes effective guide travel from configured rod length and the aftmost modeled launch-lug position. With no lugs it assumes a full-length tower launcher. The inspected loop specifically checks \class{LaunchLug}; the presence of a \class{RailButton} is not equivalent to that lug-length correction. The rod model constrains motion but does not simulate rail bending, button contact friction, or mechanical release loads.

After burnout, aerodynamic forces, gravity, wind, and changing configuration continue to govern coast. The engine detects apogee and altitude transitions, schedules recovery/staging rules, and records warnings and milestones. If deployment occurs before the natural coast apex, an auxiliary optimum-coast calculation can estimate the altitude and time that would have been reached without that deployment. This supports optimum-delay reporting; it is not a second observed flight history.

\section{What a numerical step means}
A Runge--Kutta step forms provisional positions, velocities, orientations, and rates at prescribed substep times, evaluates the derivatives there, and combines them into the accepted endpoint. The quaternion represents orientation without Euler-angle singularities, but it must remain normalized and used with the correct transform direction. The world angular-rate vector and body aerodynamic rates are related by rotation; they are not interchangeable arrays of three numbers.

The stepper records the force/mass/condition quantities associated with its accepted data point and supplies calculated events to the engine. Event-aware limits prevent stepping blindly beyond a known ignition, ejection, or other scheduled boundary. The retained RK4 and MIT RK6 implementations have different stage formulas and can have different limiting behavior; their names alone do not establish a formal adaptive local-error tolerance. The ordinary maximum timestep and angular-step settings express the actual available controls.

Exceptions separate cancellation, calculation failure, listener termination, and other run failures. System listeners such as ground-hit or recovery-stop listeners can serve special computations with different stopping rules. A headless caller must still interpret warnings/events and close run resources, rather than merely checking that the method returned an object.

\section{Tumble and grounded modes}
The tumble stepper is a translational empirical descent model. It combines symmetric-body projected area with fin planform area and fin-count efficiency factors. The coefficients in \class{BasicTumbleStepper} are 0.56 for the body contribution and 1.42 for the effective fin contribution, normalized to rocket reference area. The fin-count table approximates exposure during tumble. This is not rotational integration of a tumbling rigid body.

Landing and tumble use \class{AbstractEulerStepper}, while \class{GroundStepper} handles grounded continuation. The engine can change the stepper while retaining the branch's identity and event history. A held quaternion after deployment/tumble reflects this numerical model, not measured canopy alignment. These mode changes must be respected by plots, physical sensors, and extensions that continue through descent.

\sourceline{\source{core/src/main/java/info/openrocket/core/simulation/SimulationStatus.java}{initialization and guide length}; \source{core/src/main/java/info/openrocket/core/simulation/BasicEventSimulationEngine.java}{inherited event lifecycle}; \source{core/src/main/java/info/openrocket/core/simulation/BasicTumbleStepper.java}{BasicTumbleStepper}; \source{core/src/main/java/info/openrocket/core/simulation/GroundStepper.java}{GroundStepper}; \source{core/src/main/java/info/openrocket/core/simulation/listeners/system/OptimumCoastListener.java}{OptimumCoastListener}.}

